% ECONOMICS_PROBLEM: {"id": "I18-236", "title": "Incentives for pandemic preparedness and vaccine innovation", "jel_code": "I18", "source_id": "OP-0060"}
% !TeX program = xelatex
\documentclass[11pt,a4paper]{article}
\usepackage[margin=23mm]{geometry}
\usepackage{fontspec}
\setmainfont{Latin Modern Roman}
\usepackage{amsmath,amssymb,mathtools}
\usepackage{xcolor,enumitem,booktabs,longtable,array,xurl,hyperref}
\definecolor{muted}{HTML}{53636C}
\hypersetup{colorlinks=true,urlcolor=blue,linkcolor=blue}
\setlength{\parindent}{0pt}
\setlength{\parskip}{5pt}
\newcommand{\R}{\mathbb R}
\newcommand{\E}{\mathbb E}
\newcommand{\N}{\mathbb N}
\newcommand{\ind}{\mathbf 1}
\newcommand{\Var}{\operatorname{Var}}
\newcommand{\Cov}{\operatorname{Cov}}
\newcommand{\supp}{\operatorname{supp}}
\DeclareMathOperator*{\argmin}{arg\,min}
\DeclareMathOperator*{\argmax}{arg\,max}
\newcommand{\entrymeta}[3]{{\sffamily\small\color{muted}\textbf{\texttt{#1}}\quad|\quad #2\par #3\par}}
\newcommand{\Problemref}[1]{\texttt{#1}}
\newcommand{\JELref}[2]{\texttt{#2} (JEL \texttt{#1})}
\begin{document}
\section*{Incentives for pandemic preparedness and vaccine innovation}
\textbf{Problem ID:} \texttt{I18-236}\quad \textbf{JEL:} \texttt{I18}\par
% BEGIN ECONOMICS STATEMENT
\label{problem:OP-0060}
\label{atlas:prob:H10}
\noindent\textbf{Original atlas identifier:} H10; entry 60. \textbf{Field:} Health economics. \textbf{Original tractability label:} Hard.

\noindent\textbf{Classification:} Parameterized theoretical/empirical research agenda; not a certified unsolved theorem.

\subsection*{Definitions and assumptions}
Let $\zeta=(\omega,\boldsymbol\theta,\xi)$ collect pathogen state $\omega\in\Omega$, the vector $\boldsymbol\theta$ of firms' private development/capacity cost types, and all other exogenous clinical/production shocks $\xi$. Specify their joint law $\mu$ or a declared ambiguity set $\mathcal U$ of joint laws; independence is not assumed. Firms privately observe development/capacity cost type $\theta$ and choose research $r$, maintained readiness $k$, and delivery effort after an outbreak. A mechanism $p$ specifies upfront support, readiness payments, advance purchases, eligibility, delivery milestones, and state-contingent rewards. A model $m$ maps these actions into discovery, production speed, safety/approval, stockpile decay, and epidemic outcomes, with market equilibrium selection. Let $H_j(p,\zeta;m)$ denote country $j$'s monetized health and economic losses, $C(p,\zeta;m)$ real research/readiness/delivery costs, $F(p,\zeta;m)$ real financing distortions not included in $H_j$ or $C$, and $w_j\geq0$ global welfare weights. Fiscal commitments have a stated budget and financing rule; count each real loss once. Countries are indexed by $j=1,\ldots,J$ for fixed finite $J$; all losses and costs have finite expectations under every admissible joint law.

\subsection*{Formal research question}
\emph{Editorial formulation: the mathematical target below is not a theorem quoted from the references.}

Characterize mechanisms maximizing expected net benefit
\[
 \sup_{p\in\mathcal P}\mathbb E_{\zeta\sim\mu}\left[-\sum_{j=1}^{J}w_jH_j(p,\zeta;m)-C(p,\zeta;m)-F(p,\zeta;m)\right],
\]
or, if ambiguity is part of the stated objective, replace the expectation by $\inf_{\mu\in\mathcal U}\mathbb E_\mu[\cdot]$. The feasible set $\mathcal P$ imposes supplier incentive compatibility, participation, contract verifiability, clinical approval, fiscal limits, and a credible commitment technology. Determine which observed procurement and development variation identifies effort, cost distributions, and readiness depreciation. Compare standing advance-market commitments, subscriptions, and stockpiles under the same risk distribution, financing rule, and distributional objective; derive $\varepsilon$-optimal contracts and sensitivity bounds rather than assuming one instrument universally dominates.

\subsection*{Interpretation and scope}
Purchase payments are transfers within a specified global boundary; real capacity, scientific effort, and financing distortions are costs. Readiness and innovation are distinct: a successful vaccine candidate may still be delivered slowly when factories or inputs are unavailable. Correlated platform failures make portfolio redundancy different from independent success probabilities and must enter the production model. A standing contract also requires commitment across political terms; assuming guaranteed future procurement without explaining enforcement omits the central incentive problem. Pandemic probabilities and benefits cannot generally be learned from a single crisis, so transparent scenario or ambiguity analysis is essential. The substantive extension concerns incentives before a specific epidemic and joint readiness/innovation, while existing preventive-treatment and outbreak pull-funding mechanisms are acknowledged solved models.

\subsection*{Source audit and status}
The two named works are verified. The third original item, “advance-market-commitment literature,” is not a unique citation and remains flagged in [S3]. [S4] is an explicit editorial reference to Kremer--Levin--Snyder (2022), with its 2020 working-paper antecedent. These works do not certify a literature-wide unresolved standing-contract theorem in 2026.

\subsection*{References}
\begin{enumerate}[label={[S\arabic*]},leftmargin=*,itemsep=0.6em]
\item Michael Kremer and Christopher M. Snyder (2015). Preventives versus Treatments. The Quarterly Journal of Economics 130(3), 1167--1239. DOI: 10.1093/qje/qjv012.
\url{https://academic.oup.com/qje/article-abstract/130/3/1167/1931641}
\newline\emph{Locator and support limits:} Abstract and model: demand heterogeneity and incentives for prevention versus treatment; no standing preparedness-contract solution.
\item Christopher M. Snyder, Kendall Hoyt, Dimitrios Gouglas, Thomas Johnston, and James Robinson (2020). Designing Pull Funding for a COVID-19 Vaccine. Health Affairs 39(9), 1633--1642. DOI: 10.1377/hlthaff.2020.00646.
\url{https://www.healthaffairs.org/doi/full/10.1377/hlthaff.2020.00646}
\newline\emph{Locator and support limits:} Abstract and mechanism: pull funding during COVID-19; model/simulation recommendations are not standing-policy causal estimates.
\item Unresolved original: “advance-market-commitment literature” names a research area, not a unique bibliographic source.
\url{https://www.nber.org/papers/w28168}
\newline\emph{Locator and support limits:} A specifically identified AMC source is supplied separately; it is an editorial addition.
\item Michael Kremer, Jonathan D. Levin, and Christopher M. Snyder (2022). Designing Advance Market Commitments for New Vaccines. Management Science 68(7), 4786--4814. Earlier version: NBER Working Paper 28168 (2020), DOI 10.3386/w28168.
\url{https://www.nber.org/papers/w28168}
\newline\emph{Locator and support limits:} Abstract and model: investment hold-up, cost uncertainty, and competition; not the exact original literature-label citation.
\end{enumerate}

\subsection*{Common conventions: Source mathematical conventions}

Unless an entry states otherwise, agent and firm sets are finite. State and action spaces are measurable, strategies and outcome maps are measurable, probability laws are specified, and expectations exist for the targets under discussion. Infinite-horizon discount factors lie in $(0,1)$ and discounted objectives are finite. Norms, tolerances and complexity input sizes must be specified for computational comparisons. Constants or primitives that appear locally are inputs, not empirical facts asserted by this catalogue.

\textbf{Identification.} A statistical model $m\in\mathcal M$ implies an observable law $P_m$ and target $\theta(m)$. At law $P$, its identified set is
\[
\Theta(P;\mathcal M)=\{\theta(m):m\in\mathcal M,\ P_m=P\}.
\]
Point identification holds when this set is a singleton, or equivalently when $P_{m_1}=P_{m_2}$ implies $\theta(m_1)=\theta(m_2)$ for all compatible models. Sharp bounds mean bounds attained, or approached when the boundary is not attained, by compatible models. Writing this set is a definition, not a solution: the substantive task is to establish informative restrictions and derive or estimate the set. An unrestricted-confounding model may give only trivial bounds.

\textbf{Inference.} Identification, estimability and valid uncertainty quantification are separate questions. Fix $\alpha\in(0,1)$. A confidence procedure $C_n$ over a declared law class $\mathcal P$ has asymptotic uniform coverage at level $1-\alpha$ when
\[
\liminf_{n\to\infty}\inf_{P\in\mathcal P}\Pr_P\{\theta(P)\in C_n\}\geq1-\alpha.
\]
For set-identified targets the coverage event must be defined separately, for example coverage of the full identified set. If an entry uses identification-indexed models instead of uniquely indexed laws, the same coverage convention is applied over those models. Pointwise limits do not establish uniform validity, and asymptotic validity does not guarantee finite-sample precision. Sample sizes, dependence structures and weak-identification sequences are part of each application.

\textbf{Counterfactuals.} $Y_i(a)$ is a potential outcome under a specified intervention, including its timing, implementation and versions. It is not $Y_i$ conditional on voluntarily choosing $a$. A full intervention vector permits interference; shorthand scalar treatment notation does not silently impose no interference. Assignment, selection, attrition, anticipation and measurement must be specified. A claim about an instrument's local effect is not automatically a population policy effect.

\textbf{Economic equilibrium.} A counterfactual model includes best responses, constraints, feasibility, market clearing, state transitions and expectations consistent with the stated information. When several equilibria exist, either a declared selector is an input or the target is set-valued. Comparisons across policy regimes must use the stated selection convention. Reduced-form relationships are not presumed invariant after an equilibrium-changing intervention.

\textbf{Optimization and welfare.} Optimization is over the stated feasible set. If attainment is not established, $\sup$ or $\inf$ and $\varepsilon$-optimal choices replace an asserted maximizer or minimizer. For example, with value $V=\sup_{a\in\mathcal A}W(a)<\infty$, an $\varepsilon$-optimal set is $\{a\in\mathcal A:W(a)\geq V-\varepsilon\}$ for $\varepsilon>0$. Benefits, costs, risk and health or environmental outcomes can be aggregated only using declared units and conversion weights. Interpersonal utility weights, the treatment of fiscal transfers and foreign profits, discounting, and the behavioral welfare criterion are explicit inputs. A comparative-static derivative additionally requires existence and appropriate differentiability of the selected counterfactual.

\textbf{Sources.} Local labels [S1], [S2], and so forth restart in every question. Locators identify the relevant abstract, model, section or result at the level actually checked; a bibliographic or abstract-level verification is not represented as inspection of an entire proof. Working papers are distinguished from later publications and changing versions. Source summaries are paraphrased. All URL checks and editorial formulations refer to the audit date stated above.
% END ECONOMICS STATEMENT
\end{document}
